\documentclass[UTF8,12pt,a4paper,fontset=none]{ctexart}
\usepackage[a4paper,left=1.82cm,right=1.82cm,top=1.72cm,bottom=1.82cm,headheight=15pt]{geometry}
\usepackage{fontspec,xeCJK}
\setmainfont{Liberation Serif}
\setsansfont{DejaVu Sans}
\setmonofont{DejaVu Sans Mono}
\setCJKmainfont[BoldFont=Noto Serif CJK SC Bold]{Noto Serif CJK SC}
\setCJKsansfont[BoldFont=Noto Sans CJK SC Bold]{Noto Sans CJK SC}
\setCJKmonofont{Noto Sans Mono CJK SC}
\usepackage{amsmath,amssymb,mathtools,bm}
\usepackage{booktabs,longtable,tabularx,array,multirow}
\usepackage{enumitem,xcolor}
\usepackage[most]{tcolorbox}
\usepackage{fancyhdr,titlesec,hyperref,cleveref,lastpage}
\usepackage{tikz,float,caption,listings}
\usetikzlibrary{arrows.meta,positioning,calc,fit,backgrounds,shapes.geometric}
\definecolor{mainblue}{RGB}{39,82,138}
\definecolor{deepblue}{RGB}{25,55,95}
\definecolor{softblue}{RGB}{235,243,252}
\definecolor{softgreen}{RGB}{238,248,241}
\definecolor{softorange}{RGB}{255,246,232}
\definecolor{softgray}{RGB}{245,246,248}
\definecolor{warnred}{RGB}{160,48,48}
\hypersetup{unicode=true,colorlinks=true,linkcolor=deepblue,urlcolor=mainblue,citecolor=deepblue}
\setlength{\parindent}{2em}
\setlength{\parskip}{0.36em}
\setlength{\tabcolsep}{5pt}
\renewcommand{\arraystretch}{1.2}
\allowdisplaybreaks[4]
\emergencystretch=3em
\sloppy
\pagestyle{fancy}\fancyhf{}\fancyfoot[C]{\small 第 \thepage\ 页，共 \pageref{LastPage} 页}\renewcommand{\headrulewidth}{0.4pt}
\titleformat{\section}{\Large\bfseries\color{deepblue}}{\thesection}{0.65em}{}
\titleformat{\subsection}{\large\bfseries\color{mainblue}}{\thesubsection}{0.55em}{}
\titleformat{\subsubsection}{\normalsize\bfseries}{\thesubsubsection}{0.5em}{}
\numberwithin{equation}{section}
\newcommand{\R}{\mathbb{R}}
\newcommand{\E}{\mathbb{E}}
\newcommand{\Prob}{\mathbb{P}}
\newcommand{\Loss}{\mathcal{L}}
\newcommand{\norm}[1]{\left\lVert #1\right\rVert}
\newcommand{\ip}[2]{\left\langle #1,#2\right\rangle}
\newcommand{\tr}{\operatorname{Tr}}
\newcommand{\diag}{\operatorname{diag}}
\newcommand{\rank}{\operatorname{rank}}
\newcommand{\softmax}{\operatorname{softmax}}
\newcommand{\argmin}{\operatorname*{arg\,min}}
\newcommand{\argmax}{\operatorname*{arg\,max}}
\newcommand{\sg}{\operatorname{sg}}
\newcommand{\KL}{D_{\mathrm{KL}}}
\newcommand{\dd}{\mathrm{d}}
\newtcolorbox{keybox}[1][]{enhanced,breakable,colback=softblue,colframe=mainblue,boxrule=0.8pt,arc=1.5mm,left=2mm,right=2mm,top=1.2mm,bottom=1.2mm,title={#1},fonttitle=\bfseries}
\newtcolorbox{examplebox}[1][]{enhanced,breakable,colback=softgreen,colframe=green!45!black,boxrule=0.7pt,arc=1.5mm,left=2mm,right=2mm,top=1.2mm,bottom=1.2mm,title={#1},fonttitle=\bfseries}
\newtcolorbox{notebox}[1][]{enhanced,breakable,colback=softorange,colframe=orange!70!black,boxrule=0.7pt,arc=1.5mm,left=2mm,right=2mm,top=1.2mm,bottom=1.2mm,title={#1},fonttitle=\bfseries}
\newtcolorbox{criticalbox}[1][]{enhanced,breakable,colback=red!3,colframe=warnred,boxrule=0.8pt,arc=1.5mm,left=2mm,right=2mm,top=1.2mm,bottom=1.2mm,title={#1},fonttitle=\bfseries}
\lstset{basicstyle=\ttfamily\small,breaklines=true,frame=single,backgroundcolor=\color{softgray},numbers=left,numberstyle=\tiny,showstringspaces=false,columns=fullflexible}
\hypersetup{pdftitle={12 MUDDFormer 数学过程详解},pdfauthor={OpenAI}}
\fancyhead[L]{\small 12 MUDDFormer}
\fancyhead[R]{\small 数学过程详解}
\setlength{\abovedisplayskip}{0.82em}
\setlength{\belowdisplayskip}{0.82em}
\setlength{\abovedisplayshortskip}{0.45em}
\setlength{\belowdisplayshortskip}{0.45em}
\begin{document}
\begin{center}
{\LARGE\bfseries 12\_MUDDFormer 数学过程详解}\par
\vspace{0.35em}
{\large 多路动态稠密连接、深度注意力类比与复杂度推导}
\end{center}
\vspace{0.45em}
\begin{keybox}[本文只处理真正需要推导的部分]
本文聚焦动态跨层权重、Q/K/V/R 四路分解、深度方向注意力类比、参数量与 FLOPs 的逐步化简。全文按“公式从哪里来--每个符号是什么--中间步骤为何成立--维度和复杂度如何核对--论文结论依赖哪些假设”的顺序展开；不复述实验表格，也不使用与论文无关的通用背景凑页数。
\end{keybox}
\tableofcontents
\newpage

\section{从普通残差到跨层稠密连接}
标准 Transformer 第 $i$ 层只接收上一层表示 $X_{i-1}\in\R^{T\times D}$。MUDDFormer 保留全部历史表示
\begin{equation}
 X_{:i}=(X_0,X_1,\ldots,X_i),
\end{equation}
并在下一层输入前进行跨层加权。若使用静态权重 $a_{ij}$，则
\begin{equation}
 \bar X_i=\sum_{j=0}^{i}a_{ij}X_j.
\end{equation}
普通残差连接是特例：只令 $a_{i,i}=1$，其余为零；DenseFormer 则让 $a_{ij}$ 成为与 token 无关的可学习标量。

\section{动态稠密聚合的张量公式}
MUDD 的权重依赖当前 token 表示。对第 $i$ 层，权重生成器为
\begin{equation}
 A_i(X_i)=\operatorname{GELU}(\operatorname{RMSNorm}(X_i)W_1)W_2+a_i.
\end{equation}
维度为
\begin{align}
 X_i&\in\R^{T\times D},\\
 W_1&\in\R^{D\times K_i},\\
 W_2&\in\R^{K_i\times(i+1)},\\
 a_i&\in\R^{i+1},\\
 A_i&\in\R^{T\times(i+1)}.
\end{align}
第 $t$ 个 token 对不同深度的权重是 $A_i[t,:]$。最终聚合
\begin{equation}
 \boxed{\bar X_i[t,:]=\sum_{j=0}^{i}A_{ij}[t]X_j[t,:].}
\end{equation}
每个 token 可以选择不同深度来源，因此这是“位置相关、深度方向”的动态连接。

\section{RMSNorm 与权重尺度}
对向量 $x\in\R^D$，RMSNorm 为
\begin{equation}
 \operatorname{RMSNorm}(x)=g\odot\frac{x}{\sqrt{D^{-1}\sum_{d=1}^{D}x_d^2+\epsilon}}.
\end{equation}
若不归一化，$X_i$ 随深度幅值增长会直接放大 $W_1$ 输入，使连接权重分布随层数漂移。RMSNorm 只控制二阶尺度，不减去均值，因此比 LayerNorm 少一个中心化步骤。

\section{为什么可视为深度方向注意力}
普通自注意力在位置 $t$ 的形式为
\begin{equation}
 \operatorname{Attn}(x_t,x_{:t})
 =\softmax\left(x_tW^Q(x_{:t}W^K)^\top\right)x_{:t}W^V.
\end{equation}
MUDD 在深度方向的对应式为
\begin{equation}
 \operatorname{DA}(X_i[t],X_{:i}[t])
 =\left(\operatorname{GELU}(X_i[t]W_1)W_2+a_i\right)X_{:i}[t].
\end{equation}
对应关系为：
\begin{center}
\begin{tabular}{lll}
\toprule
横向注意力 & 深度动态连接 & 差异\\
\midrule
查询 $x_tW^Q$ & $\operatorname{GELU}(X_i[t]W_1)$ & 查询由当前层 token 产生\\
键 $x_{:t}W^K$ & 参数矩阵 $W_2^\top$ & 键与输入无关\\
Softmax & GELU 后线性输出 & 权重不归一化\\
值投影 $W^V$ & 原始 $X_j[t]$ & 无额外值投影\\
\bottomrule
\end{tabular}
\end{center}
所以它不是标准注意力的严格等价物，而是去掉 Softmax、输入依赖键和值投影后的深度线性注意力近似。

\section{不使用 Softmax 的后果}
若使用 Softmax，则
\begin{equation}
 \sum_{j=0}^{i}A_{ij}[t]=1,
\end{equation}
输出是历史表示的凸组合。MUDD 不做归一化，所以权重可为负，且总和不受限：
\begin{equation}
 \bar X_i[t]=\sum_j A_{ij}[t]X_j[t].
\end{equation}
这使模型可以执行差分、放大与抵消，但也可能导致尺度不稳定。论文通过 RMSNorm、静态先验 $a_i$ 和残差结构控制该问题。

\section{Multiway：Q、K、V、R 四条深度路径}
普通 Transformer 同一个输入同时生成查询、键、值并参与残差。MUDDFormer 解耦四路：
\begin{equation}
 (X_i^Q,X_i^K,X_i^V,X_i^R)
 =\left(\operatorname{DA}_i^Q(X_{:i}),
 \operatorname{DA}_i^K(X_{:i}),
 \operatorname{DA}_i^V(X_{:i}),
 \operatorname{DA}_i^R(X_{:i})\right).
\end{equation}
随后
\begin{align}
 X_A'&=\operatorname{MHA}(\operatorname{LN}X_i^Q,
 \operatorname{LN}X_i^K,\operatorname{LN}X_i^V)+X_i^R,\\
 X_i&=\operatorname{FFN}(\operatorname{LN}X_A')+X_A'.
\end{align}
四路连接权重不必相同。查询可能偏好高层语义，键可能保留低层局部模式，值和残差又可选择不同历史层。

\section{动态权重生成器的参数量}
第 $i$ 层四路连接需要输出 $4(i+1)$ 个权重，因此隐藏宽度记为
\begin{equation}
 K_i=4(i+1).
\end{equation}
主要参数为
\begin{equation}
 P_i=DK_i+K_i^2+K_i,
\end{equation}
分别来自 $W_1,W_2,a_i$。忽略偏置并对所有层求和：
\begin{equation}
 P_{\rm extra}=\sum_{i=1}^{L}(DK_i+K_i^2).
\end{equation}
标准 Transformer 每层注意力约 $4D^2$、FFN 约 $8D^2$，总参数近似
\begin{equation}
 P_{\rm base}=12LD^2.
\end{equation}
所以
\begin{equation}
 R_{\Delta P}=\frac{\sum_i(DK_i+K_i^2)}{12LD^2}.
\end{equation}

\section{平均层近似与论文的 $\eta/6$ 公式}
取中间层平均隐藏维度
\begin{equation}
 K\approx2(L+3).
\end{equation}
当 $D\gg K$ 时忽略 $K^2$：
\begin{align}
 R_{\Delta P}
 &\approx\frac{DK}{12D^2}\\
 &=\frac{K}{12D}\\
 &=\frac{2(L+3)}{12D}\\
 &=\boxed{\frac{L+3}{6D}}.
\end{align}
定义深宽比
\begin{equation}
 \eta=\frac{L+3}{D},
\end{equation}
则
\begin{equation}
 R_{\Delta P}\approx\frac{\eta}{6}.
\end{equation}
这说明额外参数主要由“深度相对宽度”决定。宽模型中 $D$ 大，开销比例更小；极深窄模型中比例上升。

\section{额外 FLOPs 的推导}
第 $i$ 层生成连接权重需要
\begin{equation}
 2TDK_i+2TK_i^2
\end{equation}
次乘加，跨层聚合再需要约
\begin{equation}
 2TDK_i.
\end{equation}
故
\begin{equation}
 F_{\rm extra}=\sum_i(4TDK_i+2TK_i^2).
\end{equation}
标准 Transformer 预训练 FLOPs 近似
\begin{equation}
 F_{\rm base}=2LDT(12D+T).
\end{equation}
平均层近似并忽略 $K^2$：
\begin{align}
 R_{\Delta F}
 &\approx\frac{4TDK}{2DT(12D+T)}\\
 &=\frac{2K}{12D+T}\\
 &=\frac{4(L+3)}{12D+T}.
\end{align}
令 $\rho=T/D$，则
\begin{equation}
 \boxed{R_{\Delta F}\approx\frac{\eta}{3+\rho/4}}.
\end{equation}
序列越长，基础注意力成本增大，MUDD 的相对开销反而下降。

\section{梯度如何跨越深度}
若
\begin{equation}
 \bar X_i=\sum_{j=0}^{i}A_{ij}X_j,
\end{equation}
则对历史层 $X_j$ 的梯度有两部分：
\begin{equation}
 \frac{\partial\Loss}{\partial X_j}
 =A_{ij}\frac{\partial\Loss}{\partial\bar X_i}
 +\frac{\partial A_i}{\partial X_j}^{\top}
 \left(X_{:i}\frac{\partial\Loss}{\partial\bar X_i}\right).
\end{equation}
因为 $A_i$ 只由当前 $X_i$ 生成，$j<i$ 时第二项通常为零，历史层获得一条由 $A_{ij}$ 加权的直接梯度路径。这是缓解深层残差瓶颈的核心机制。

\section{极限情形核对}
\begin{itemize}[leftmargin=2em]
\item 若 $A_{ii}=1$、其余为零，退化为普通逐层 Transformer；
\item 若 $A_{ij}$ 仅是静态参数，退化为静态 DenseFormer；
\item 若四路 $Q,K,V,R$ 共享同一 $A_i$，退化为单路动态稠密连接；
\item 若再对 $A_i$ 做 Softmax，则输出变成深度方向凸组合，表达能力与稳定性折中改变。
\end{itemize}

\section{容易误解的地方}
\begin{criticalbox}[“深度注意力”只是结构类比]
MUDD 的键是参数而不是输入，权重没有 Softmax，值没有独立投影，因此不能把它当作标准自注意力直接套用概率解释。论文的参数/FLOPs 公式使用平均层与 $D\gg K$ 近似；在极深窄模型上，$K_i^2$ 项不可忽略。
\end{criticalbox}
\section{最终公式路线图}
\begin{keybox}[阅读完成后应掌握]
本文的数学主线已经被压缩为：先明确随机变量或张量的定义，再由概率分解、微分方程、优化目标或矩阵运算得到论文核心公式；最后用维度、极限情形、参数量和复杂度进行反向核验。遇到论文中的简写式，应先恢复被省略的条件变量、求和指标与归一化常数，再讨论其算法含义。
\end{keybox}
\clearpage
\section{动态稠密连接的统一矩阵写法}
设 Transformer 有 $L$ 层，隐藏状态矩阵
\begin{equation}
 H=[h_0,h_1,\ldots,h_{L-1}]\in\R^{d\times L}.
\end{equation}
第 $\ell$ 层不再只读取 $h_{\ell-1}$，而是生成深度权重
\begin{equation}
 a_\ell=g_\ell(h_{0:\ell-1})\in\R^{\ell},
\end{equation}
并聚合
\begin{equation}
 \widetilde h_\ell=H_{<\ell}a_\ell
 =\sum_{j=0}^{\ell-1}a_{\ell j}h_j.
\end{equation}
若 Q、K、V、残差分别使用独立权重，记为 $a_\ell^Q,a_\ell^K,a_\ell^V,a_\ell^R$，则
\begin{equation}
 q_\ell=W_QH_{<\ell}a_\ell^Q,
 \quad
 k_\ell=W_KH_{<\ell}a_\ell^K,
 \quad
 v_\ell=W_VH_{<\ell}a_\ell^V,
\end{equation}
\begin{equation}
 r_\ell=H_{<\ell}a_\ell^R.
\end{equation}
这就是 Multiway 的精确线性代数含义：同一组历史层状态，经过四个不同的深度方向选择器。

\section{深度注意力与普通 token 注意力的对应关系}
普通注意力对 token 维做
\begin{equation}
 \operatorname{Attn}(Q,K,V)=\softmax\left(\frac{QK^\top}{\sqrt d}\right)V.
\end{equation}
动态稠密连接则把“键值集合”换成历史层：
\begin{equation}
 a_{\ell j}=\frac{\exp(s_{\ell j})}{\sum_{k<\ell}\exp(s_{\ell k})},
 \qquad
 \widetilde h_\ell=\sum_{j<\ell}a_{\ell j}h_j.
\end{equation}
如果论文不使用 softmax，而直接输出实值权重，则它不再是凸组合，可能出现
\begin{equation}
 \sum_j a_{\ell j}\ne1,
 \qquad a_{\ell j}<0.
\end{equation}
优点是能做差分和放大；缺点是范数没有天然上界。

\section{RMSNorm 对深度聚合尺度的作用}
RMSNorm 定义为
\begin{equation}
 \operatorname{RMSNorm}(h)=
 \gamma\odot\frac{h}{\sqrt{d^{-1}\|h\|_2^2+\epsilon}}.
\end{equation}
忽略 $\epsilon$ 和逐维 $\gamma$ 时，归一化后
\begin{equation}
 \frac1d\|\operatorname{RMSNorm}(h)\|_2^2\approx1.
\end{equation}
若历史层范数随深度增长，直接计算权重可能把“尺度大”误判为“更重要”。先归一化再生成 $a_{\ell j}$，可以让权重更多反映方向和内容，而不是绝对幅值。

\section{三层玩具例子的完整前向计算}
设标量隐藏状态 $h_0=1,h_1=2,h_2=-1$。第 3 层的四组权重为
\begin{equation}
 a^Q=(0.2,0.5,0.3),\quad
 a^K=(0.1,0.1,0.8),
\end{equation}
\begin{equation}
 a^V=(0.6,0.2,0.2),\quad
 a^R=(0.3,0.4,0.3).
\end{equation}
则
\begin{align}
 \widetilde h^Q&=0.2(1)+0.5(2)+0.3(-1)=0.9,\\
 \widetilde h^K&=0.1(1)+0.1(2)+0.8(-1)=-0.5,\\
 \widetilde h^V&=0.6(1)+0.2(2)+0.2(-1)=0.8,\\
 r&=0.3(1)+0.4(2)+0.3(-1)=0.8.
\end{align}
即使来自同一历史，Q/K/V/R 看到的有效深度表示也不同。这是比单一 DenseNet 式求和更高的自由度。

\section{梯度图是块下三角结构}
第 $\ell$ 层依赖所有 $j<\ell$ 的隐藏状态，因此总 Jacobian
\begin{equation}
 J_{\ell j}=\frac{\partial h_\ell}{\partial h_j}
\end{equation}
在层索引上是严格下三角加对角结构。若
\begin{equation}
 h_\ell=F_\ell(\widetilde h_\ell)+r_\ell,
\end{equation}
则
\begin{equation}
 \frac{\partial h_\ell}{\partial h_j}
 =J_F\left(a_{\ell j}I+
 \sum_k h_k\frac{\partial a_{\ell k}}{\partial h_j}\right)
 +a^R_{\ell j}I+
 \sum_k h_k\frac{\partial a^R_{\ell k}}{\partial h_j}.
\end{equation}
第一类项是固定权重的直接路径，第二类项来自动态权重对历史状态的依赖。后者允许内容依赖的信用分配，但也引入额外二阶耦合。

\section{谱范数稳定条件}
若忽略权重生成器的导数，聚合矩阵为
\begin{equation}
 \widetilde H=HA.
\end{equation}
则
\begin{equation}
 \|\widetilde H\|_F\le\|H\|_F\|A\|_2.
\end{equation}
若每列 $a_\ell$ 是非负且和为 1，$\|A\|_1=1$，但 $\|A\|_2$ 仍可能随层数增长。无 softmax 时更需要显式正则，例如
\begin{equation}
 \lambda\sum_\ell\|a_\ell\|_2^2
\end{equation}
或门控缩放，避免历史层线性组合无限放大。

\section{参数量与 FLOPs 的逐项核对}
若每层用一个两层 MLP 从 $d$ 维状态生成长度 $m$ 的深度权重，瓶颈维度为 $r$，参数量约
\begin{equation}
 P_g=dr+rm+r+m.
\end{equation}
四路权重若完全独立，约为 $4P_g$；若共享第一层，仅分开输出头，则
\begin{equation}
 P_g^{\rm shared}=dr+r+4(rm+m).
\end{equation}
聚合本身对每个 token、每个通道需要 $m$ 次乘加，FLOPs 约
\begin{equation}
 C_{\rm agg}=2BTdm
\end{equation}
（乘和加各计一次）。当平均可见历史层数为 $\bar m$，总额外开销近似
\begin{equation}
 C_{\rm total}\approx 4\cdot2BTdL\bar m+C_g.
\end{equation}
因此论文报告的低百分比开销依赖于 $\bar m$ 的截断、权重生成器的小尺寸以及主干 FFN 的高基线 FLOPs。

\section{平均层近似的误差}
若用最近 $m$ 层均值替代动态权重，
\begin{equation}
 \bar h=\frac1m\sum_{j=1}^{m}h_j.
\end{equation}
动态聚合与均值之差
\begin{equation}
 \widetilde h-\bar h=
 \sum_{j=1}^{m}\left(a_j-\frac1m\right)h_j.
\end{equation}
由 Cauchy--Schwarz，
\begin{equation}
 \|\widetilde h-\bar h\|_2
 \le\left\|a-\frac1m\mathbf 1\right\|_2
 \left(\sum_j\|h_j\|_2^2\right)^{1/2}.
\end{equation}
当动态权重接近均匀分布时，简单平均近似良好；若权重尖锐地选择某个远层，均值会严重稀释该层信息。
\end{document}
